\pagestyle{empty}
\tight
\begin{tblr}{width=\textwidth,colspec={|Q[c,m,1.9cm]|X[l,m]|},hlines,vlines,hspan=minimal,row{1}={bg=headblue},rowsep=1pt,colsep=3pt}
\SetCell{c,m}\hfil\logo\hfil & \textbf{POLITEKNIK NEGERI MALANG}\par\textbf{JURUSAN TEKNOLOGI INFORMASI}\par\textbf{PROGRAM STUDI: D4 TEKNIK INFORMATIKA} \\
\SetCell[c=2]{c} \textbf{RENCANA PEMBELAJARAN SEMESTER (RPS) --- DRAF KONSEP B: PENGABDIAN \& SIKAP} & \\
\end{tblr}
\begin{longtblr}{width=\textwidth,colspec={|Q[l,m,1.9cm]|X[0.85,l,m]|X[1.45,l,m]|X[0.75,l,m]|X[1.15,l,m]|},hlines,vlines,hspan=minimal,rowsep=1pt,colsep=2pt,row{1,3}={bg=headgray,font=\bfseries}}
MATA KULIAH & KODE & BOBOT (SKS)/jam & SEMESTER & TGL. PENYUSUNAN \\
Kuliah Kerja Nyata Tematik & RTI257002 (sementara)\footnotemark[1] & 2 SKS / 4 jam/minggu & 6/7 (sementara)\footnotemark[1] & 13 Juli 2026 \\
OTORISASI & Dosen Pengembang RPS & Koordinator RMK & \SetCell[c=2]{l} Ka PRODI &  \\
 & Tim Pengajar Program Studi D4 TI &  & \SetCell[c=2]{l}{\vspace{8mm}\par Dr. Ely Setyo Astuti, S.T., M.T.} &  \\
\SetCell[r=14]{l} Capaian Pembelajaran (CP) & \SetCell[c=4]{l,bg=headgray,font=\bfseries} Capaian Pembelajaran Lulusan Program Studi (CPL-Prodi) &  &  &  \\
 & CPL01 & \SetCell[c=3]{l} Menginternalisasi ketakwaan kepada Tuhan Yang Maha Esa, menaati hukum, disiplin, dan menunjukkan profesionalisme melalui etika profesi, pembelajaran sepanjang hayat, serta respons terhadap isu sosial dan teknologi. &  &  \\
 & CPL03 & \SetCell[c=3]{l} Mampu mengelola tim, manajemen diri, berkomunikasi secara lisan maupun tertulis dengan baik dan mampu melakukan presentasi. &  &  \\
 & CPL08 & \SetCell[c=3]{l} Mampu mengelola proyek teknologi informasi secara profesional dengan mengoptimalkan sumber daya yang tersedia. &  &  \\
 & \SetCell[c=4]{l,bg=headgray,font=\bfseries} Capaian Pembelajaran mata kuliah (CPL-MK) &  &  &  \\
 & CPMK0101 & \SetCell[c=3]{l} Mampu menunjukkan profesionalisme melalui penerapan etika profesi TI, keselamatan kerja, kewirausahaan bertanggung jawab, rekayasa perangkat lunak ramah lingkungan, dan disiplin dalam penelitian maupun penerapan aplikasi teknologi. &  &  \\
 & CPMK0102 & \SetCell[c=3]{l} Mampu menginternalisasi nilai ketakwaan kepada Tuhan YME, Pancasila, kewarganegaraan, dan menunjukkan pembelajaran sepanjang hayat melalui disiplin diri, tanggung jawab sosial, dan adaptasi karir profesional. &  &  \\
 & CPMK0301 & \SetCell[c=3]{l} Mampu mengelola diri secara profesional, berkomunikasi secara lisan maupun tertulis, dan melakukan presentasi secara efektif dalam konteks kerja. &  &  \\
 & CPMK0302 & \SetCell[c=3]{l} Mampu mengelola tim dan proyek secara terencana, mulai dari perencanaan, koordinasi, komunikasi antar anggota, hingga penyampaian hasil kepada pemangku kepentingan. &  &  \\
 & \SetCell[c=4]{l,bg=headgray,font=\bfseries} Kemampuan akhir tiap tahapan belajar (Sub-CPMK) &  &  &  \\
 & SCPMK0101-06001 & \SetCell[c=3]{l} Mahasiswa mampu menunjukkan etika kerja dan profesionalisme selama observasi dan pelaksanaan program pengabdian di lokasi mitra. &  &  \\
 & SCPMK0102-06002 & \SetCell[c=3]{l} Mahasiswa mampu mengidentifikasi kebutuhan sosial mitra/komunitas dan melaksanakan kegiatan pemberdayaan yang berdampak nyata. &  &  \\
 & SCPMK0301-06003 & \SetCell[c=3]{l} Mahasiswa mampu menggali kebutuhan mitra secara partisipatif dan mengomunikasikan hasil program kerja pengabdian. &  &  \\
 & SCPMK0302-06004 & \SetCell[c=3]{l} Mahasiswa mampu menyusun, mengelola, dan mengendalikan pelaksanaan program kerja pengabdian bersama tim dan mitra. &  &  \\
\SetCell[r=2]{l} Penilaian & \SetCell[c=4]{l} *Nilai CPMK didapat dari agregasi nilai SUB-CPMK yang dibebankan &  &  &  \\
 & \SetCell[c=4]{l}{\tiny\begin{tblr}{width=\linewidth,colspec={|Q[c,m,0.1\linewidth]|Q[c,m,0.16\linewidth]|X[l,m]|X[l,m]|X[l,m]|X[l,m]|Q[c,m,0.08\linewidth]|},hlines,vlines,hspan=minimal,rowsep=1pt,colsep=0.7pt,row{1,2}={bg=headgray,font=\bfseries}}
Kode CPMK & Kode SCPMK & \SetCell[c=4]{c} Bobot per Bentuk Penilaian & & & & Total Bobot \\
 & & Proposal & Logbook \& Keterlibatan & Penilaian Mitra/DPL & Laporan \& Seminar & \\
CPMK0101 & SCPMK0101-06001 & 5\%: etika \& kelayakan program & & 10\%: penilaian sikap \& profesionalisme & & 15\% \\
CPMK0102 & SCPMK0102-06002 & & 20\%: keterlibatan \& kontribusi pemberdayaan & 15\%: penilaian dampak sosial oleh mitra & & 35\% \\
CPMK0301 & SCPMK0301-06003 & 10\%: komunikasi kebutuhan dengan mitra & & & 15\%: komunikasi hasil kepada mitra dan penguji & 25\% \\
CPMK0302 & SCPMK0302-06004 & 5\%: perencanaan program kerja & 10\%: pengelolaan tim \& pelaksanaan & & 10\%: pertanggungjawaban pelaksanaan & 25\% \\
\SetCell[c=6]{r}\textbf{Total Per Penilaian} & & 20\% & 30\% & 25\% & 25\% & \textbf{100\%} \\
\end{tblr}} &  &  &  \\
Diskripsi Singkat Mata Kuliah & \SetCell[c=4]{l} Kuliah Kerja Nyata Tematik (KKN Tematik) adalah mata kuliah wajib 2 SKS penunjang IKU kegiatan mahasiswa di luar Prodi (dasar: Nota Dinas Wakil Direktur I No. 11/WADIR I/KM/2026 dan Permendiktisaintek No. 39 Tahun 2025) yang menempatkan mahasiswa pada mitra/komunitas untuk melaksanakan program pemberdayaan, edukasi, dan penguatan tanggung jawab sosial, disertai laporan dan seminar hasil. \\
Bahan Kajian / Materi Pembelajaran & \SetCell[c=4]{l}\begin{enumerate}\item Pembekalan KKN, etika, dan tanggung jawab sosial\item Observasi, sosialisasi, dan analisis kebutuhan komunitas\item Pelaksanaan program pemberdayaan dan edukasi masyarakat\item Pengelolaan tim dan komunikasi dengan mitra\item Penyusunan laporan KKN dan seminar hasil\end{enumerate} \\
\SetCell[r=2]{l} Pustaka & \SetCell[c=4]{l} \textbf{Utama:}\begin{enumerate}\item Buku Pedoman KKN Tematik Politeknik Negeri Malang.\item Panduan MBKM Kemdikbud 2020.\item Pedoman Penulisan Laporan KKN Polinema.\end{enumerate} \\
 & \SetCell[c=4]{l} \textbf{Pendukung:} Materi LMS, dokumen profil mitra/komunitas, dan bahan pemberdayaan masyarakat. \\
Media Pembelajaran & \SetCell[c=2]{l} \textbf{Software:} Aplikasi logbook digital, LMS, media sosialisasi/edukasi digital. &  & \SetCell[c=2]{l} \textbf{Hardware:} Komputer/laptop, perangkat presentasi, alat peraga kegiatan komunitas. &  \\
Nama Dosen Pengampu & \SetCell[c=4]{l} Tim Pengajar Program Studi D4 TI \\
Matakuliah Syarat & \SetCell[c=4]{l} Lulus minimal 90 SKS \\
\end{longtblr}
\begin{longtblr}{width=\textwidth,colspec={|Q[c,m,0.06\textwidth]|X[1.35,l,m]|X[1.08,l,m]|X[1.16,l,m]|X[0.7,l,m]|X[1.24,l,m]|X[1.06,l,m]|X[0.98,l,m]|Q[c,m,0.05\textwidth]|},hlines,vlines,hspan=minimal,rowhead=2,row{1,2}={bg=headgreen,font=\bfseries},rowsep=1pt,colsep=1.5pt}
Mgg.\par Ke & Kemampuan Akhir Yang Direncanakan (Sub-CP-MK) & Bahan kajian (Materi Pembelajaran) & Bentuk dan Metode Pembelajaran & Estimasi Waktu & Pengalaman Belajar Mahasiswa & Kriteria \& Bentuk Penilaian & Indikator Penilaian & Bobot Penilaian (\%) \\
(1) & (2) & (3) & (4) & (5) & (6) & (7) & (8) & (9) \\
1 & SCPMK0101-06001, SCPMK0302-06004 & Pembekalan KKN: orientasi, etika, dan kontrak belajar pengabdian. & Modalitas: Blended Learning. Bentuk: Luring/Daring. Metode: Case Method, ceramah interaktif, dan diskusi. & 1 x 2 x 50' tatap muka; persiapan mandiri. & Mahasiswa mengikuti pembekalan KKN, memahami etika pengabdian, dan menandatangani kontrak belajar. & Bentuk penilaian: Kehadiran dan partisipasi pembekalan. Mengacu rubrik RTM. & Ketepatan pemahaman prosedur; kesiapan etika pengabdian. & 2 \\
2 & SCPMK0102-06002, SCPMK0301-06003 & Observasi dan sosialisasi awal dengan mitra/komunitas. & Modalitas: Praktik Lapangan dan Bimbingan. Metode: observasi lapangan, sosialisasi, dan bimbingan dosen pembimbing lapangan (DPL). & 1 x 2 x 50' observasi lapangan; kerja mandiri. & Mahasiswa melakukan observasi dan sosialisasi awal ke mitra/komunitas untuk membangun kepercayaan dan mengenali konteks sosial. & Bentuk penilaian: Catatan observasi dan sosialisasi. Mengacu rubrik RTM. & Ketepatan pemetaan komunitas; kualitas komunikasi awal. & 2 \\
3 & SCPMK0102-06002, SCPMK0302-06004 & Analisis kebutuhan sosial dan potensi pemberdayaan komunitas. & Modalitas: Praktik Lapangan dan Bimbingan. Metode: analisis partisipatif dan bimbingan DPL. & 1 x 2 x 50' bimbingan; analisis mandiri. & Mahasiswa menganalisis kebutuhan sosial dan potensi pemberdayaan komunitas secara partisipatif bersama mitra. & Bentuk penilaian: Draft analisis kebutuhan sosial. Mengacu rubrik RTM. & Ketepatan analisis kebutuhan sosial; relevansi potensi pemberdayaan. & 3 \\
4 & SCPMK0101-06001, SCPMK0302-06004 & Penyusunan proposal program kerja pemberdayaan. & Modalitas: Blended Learning. Bentuk: Luring/Daring. Metode: Case Method, penyusunan dokumen, dan bimbingan DPL. & 1 x 2 x 50' bimbingan; penyusunan mandiri. & Mahasiswa menyusun proposal program kerja pemberdayaan yang memuat rencana kegiatan sosial, jadwal, dan target dampak. & Bentuk penilaian: Draft proposal program kerja. Mengacu rubrik RTM. & Kelengkapan proposal; kelayakan rencana kerja; ketepatan etika penyusunan. & 5 \\
5 & SCPMK0101-06001, SCPMK0301-06003, SCPMK0302-06004 & Seminar proposal program kerja. & Modalitas: Blended Learning. Bentuk: Luring/Daring. Metode: presentasi dan Case Method. & 1 x 2 x 50' seminar proposal. & Mahasiswa mempresentasikan proposal program kerja pemberdayaan di hadapan DPL dan menerima masukan untuk perbaikan. & Bentuk penilaian: Seminar proposal program kerja. Mengacu rubrik RTM. & Kelayakan program kerja; kejelasan target dampak sosial; kualitas presentasi. & 15 \\
6 & SCPMK0102-06002, SCPMK0302-06004 & Pelaksanaan program kerja tahap 1: kegiatan sosialisasi dan edukasi masyarakat. & Modalitas: Praktik Lapangan dan Bimbingan. Metode: praktik lapangan, monitoring berkala, dan bimbingan DPL. & 1 x 2 x 50' bimbingan; kerja lapangan mandiri. & Mahasiswa melaksanakan kegiatan sosialisasi dan edukasi awal kepada masyarakat sesuai program kerja yang disepakati. & Bentuk penilaian: Logbook dan dokumentasi kegiatan. Mengacu rubrik RTM. & Ketepatan waktu; kualitas pelaksanaan; keterlibatan masyarakat. & 3 \\
7 & SCPMK0102-06002 & Pelaksanaan program kerja tahap 1 (lanjutan): pendampingan komunitas. & Modalitas: Praktik Lapangan dan Bimbingan. Metode: praktik lapangan, monitoring berkala, dan bimbingan DPL. & 1 x 2 x 50' bimbingan; kerja lapangan mandiri. & Mahasiswa mendampingi komunitas secara berkelanjutan dan mendokumentasikan perkembangan kegiatan pemberdayaan. & Bentuk penilaian: Logbook dan progress pendampingan. Mengacu rubrik RTM. & Kualitas pendampingan; ketepatan waktu; keterlibatan mitra. & 3 \\
8 & SCPMK0101-06001, SCPMK0302-06004 & Monitoring dan evaluasi (monev) tengah bersama mitra dan DPL. & Modalitas: Praktik Lapangan dan Bimbingan. Metode: praktik lapangan, monitoring berkala, dan bimbingan DPL. & Disesuaikan jadwal mitra. & Mahasiswa menjalani monev tengah program KKN bersama mitra dan DPL untuk mengevaluasi kemajuan program pemberdayaan. & Bentuk penilaian: Formulir monev tengah oleh mitra/DPL. Mengacu rubrik RTM. & Kesesuaian progres dengan rencana; profesionalisme; kualitas kemajuan program. & 15 \\
9 & SCPMK0102-06002 & Pelaksanaan program kerja tahap 2: kegiatan pemberdayaan lanjutan. & Modalitas: Praktik Lapangan dan Bimbingan. Metode: praktik lapangan, monitoring berkala, dan bimbingan DPL. & 1 x 1 x 50' monitoring; kerja mandiri. & Mahasiswa melanjutkan kegiatan pemberdayaan berdasarkan masukan hasil monev tengah. & Bentuk penilaian: Logbook dan dokumentasi kegiatan lanjutan. Mengacu rubrik RTM. & Kualitas pelaksanaan; ketepatan waktu. & 3 \\
10 & SCPMK0301-06003, SCPMK0102-06002 & Pelaksanaan program kerja tahap 2 (lanjutan): pelatihan dan edukasi kepada masyarakat. & Modalitas: Praktik Lapangan dan Bimbingan. Metode: praktik lapangan, monitoring berkala, dan bimbingan DPL. & 1 x 1 x 50' monitoring; kerja mandiri. & Mahasiswa melatih dan mengedukasi masyarakat sesuai program kerja serta mengumpulkan umpan balik peserta. & Bentuk penilaian: Logbook dan berita acara pelatihan. Mengacu rubrik RTM. & Kualitas pelatihan; ketepatan waktu; keterlibatan masyarakat. & 3 \\
11 & SCPMK0302-06004 & Pelaksanaan program kerja tahap 3: evaluasi dampak sosial program. & Modalitas: Praktik Lapangan dan Bimbingan. Metode: praktik lapangan, monitoring berkala, dan bimbingan DPL. & 1 x 1 x 50' monitoring; kerja mandiri. & Mahasiswa mengevaluasi dampak sosial program kerja terhadap komunitas/mitra. & Bentuk penilaian: Logbook dan catatan evaluasi dampak sosial. Mengacu rubrik RTM. & Ketepatan waktu; kualitas evaluasi dampak. & 3 \\
12 & SCPMK0102-06002 & Konsolidasi luaran dan dokumentasi dampak sosial. & Modalitas: Praktik Lapangan dan Bimbingan. Metode: praktik lapangan, monitoring berkala, dan bimbingan DPL. & 1 x 1 x 50' monitoring; kerja mandiri. & Mahasiswa menyusun dokumentasi dampak sosial program pemberdayaan sebagai luaran kegiatan. & Bentuk penilaian: Dokumentasi dampak sosial. Mengacu rubrik RTM. & Kelengkapan dokumentasi; ketepatan waktu. & 3 \\
13 & SCPMK0101-06001, SCPMK0301-06003 & Bimbingan penyusunan laporan KKN. & Modalitas: Blended Learning. Bentuk: Luring/Daring. Metode: Case Method dan bimbingan DPL. & 1 x 1 x 50' bimbingan; penulisan mandiri. & Mahasiswa mendapatkan bimbingan teknis penulisan laporan KKN mencakup sistematika, konten, dan tata tulis ilmiah. & Bentuk penilaian: Draft laporan dengan catatan bimbingan. Mengacu rubrik RTM. & Ketepatan waktu; kualitas dokumen. & 5 \\
14 & SCPMK0301-06003 & Finalisasi laporan KKN. & Modalitas: Blended Learning. Bentuk: Luring/Daring. Metode: Case Method dan bimbingan DPL. & 1 x 1 x 50' bimbingan; penulisan mandiri. & Mahasiswa menyelesaikan draft akhir laporan KKN mencakup seluruh bab dan lampiran. & Bentuk penilaian: Laporan KKN final. Mengacu rubrik RTM. & Kelengkapan laporan; ketepatan waktu. & 5 \\
15 & SCPMK0101-06001, SCPMK0302-06004 & Persiapan seminar hasil: penyusunan bahan presentasi dan kelengkapan dokumen. & Modalitas: Blended Learning. Bentuk: Luring/Daring. Metode: Case Method dan bimbingan DPL. & 1 x 1 x 50' persiapan; mandiri. & Mahasiswa menyiapkan bahan presentasi seminar hasil dan melengkapi dokumen pendukung program kerja. & Bentuk penilaian: Kelengkapan dokumen seminar. Mengacu rubrik RTM. & Kelengkapan dokumen; kesiapan presentasi. & 5 \\
16 & SCPMK0301-06003, SCPMK0102-06002 & Seminar hasil KKN: presentasi dan defense laporan pengabdian. & Modalitas: Blended Learning. Bentuk: Luring/Daring. Metode: presentasi dan Case Method. & 1 x 2 x 50' seminar hasil. & Mahasiswa mempresentasikan dan mempertahankan laporan KKN beserta dampak sosial program di hadapan tim penguji. & Bentuk penilaian: Seminar hasil (presentasi dan tanya jawab). Mengacu rubrik RTM. & Kualitas presentasi; ketepatan jawaban; kejelasan dampak sosial. & 25 \\
\end{longtblr}
\footnotetext[1]{\textit{Draf konsep untuk perbandingan. Nama mata kuliah dan bobot 2 SKS mengikuti Nota Dinas Wakil Direktur I No. 11/WADIR I/KM/2026. Kode mata kuliah, semester, dan integrasi ke crosswalk kurikulum akan difinalkan pada versi yang dipilih.}}
